%############################################################################
%  CHAPITRE 4 — RÉALISATION DE LA PREMIÈRE RELEASE
%############################################################################
\chapter{Réalisation de la première release}
\minitoc
\vspace{0.6cm}

\section*{Introduction}
Ce chapitre est consacré à la réalisation de la première release. Nous présentons d'abord
l'environnement de travail (matériel, logiciels, langages et frameworks) ainsi que
l'architecture applicative mise en œuvre. Nous détaillons ensuite le déroulement des deux
premiers sprints, en illustrant le résultat par des captures d'écran de l'application.

\section{Environnement de travail}

\subsection{Environnement matériel}
Le développement a été réalisé sur un poste de travail disposant des caractéristiques
suivantes : processeur multi-cœurs, mémoire vive de 16 Go et système d'exploitation Windows.
Cette configuration permet d'exécuter simultanément l'ensemble des conteneurs de la plateforme
en environnement de développement.

\subsection{Environnement logiciel}
Le tableau suivant récapitule les principaux outils et technologies utilisés.

\begin{table}[H]\centering\renewcommand{\arraystretch}{1.3}
\begin{tabular}{|p{3.6cm}|p{10.2cm}|}\hline
\rowcolor{gray!15}\textbf{Technologie / Outil} & \textbf{Rôle} \\ \hline
\textbf{Angular} & Framework front-end pour la construction de l'interface utilisateur. \\ \hline
\textbf{Spring Boot} & Framework back-end pour le développement des microservices. \\ \hline
\textbf{Spring Cloud} & Composants pour l'architecture répartie (Gateway, Eureka, Config). \\ \hline
\textbf{Keycloak} & Serveur d'authentification et de gestion des identités (OAuth2 / OIDC). \\ \hline
\textbf{PostgreSQL} & Système de gestion de base de données relationnelle. \\ \hline
\textbf{MongoDB} & Base de données documentaire (NoSQL). \\ \hline
\textbf{RabbitMQ} & Courtier de messages pour la communication asynchrone. \\ \hline
\textbf{Docker} & Conteneurisation et orchestration des services (Docker Compose). \\ \hline
\textbf{OpenShift} & Plateforme cloud d'orchestration (Kubernetes) pour le déploiement. \\ \hline
\textbf{Google Gemini} & Modèle de langage utilisé par l'assistant intelligent. \\ \hline
\textbf{Maven} & Outil de construction et de gestion des dépendances Java. \\ \hline
\textbf{Git} & Système de gestion de versions du code source. \\ \hline
\textbf{Visual Studio Code / IntelliJ} & Environnements de développement intégrés. \\ \hline
\end{tabular}
\caption{Environnement logiciel du projet}\label{tab:env}
\end{table}

\subsection{Langages de programmation}
\begin{itemize}[leftmargin=1.2cm]
  \item \textbf{Java} : langage utilisé pour l'ensemble des microservices back-end.
  \item \textbf{TypeScript} : langage utilisé pour le développement front-end avec Angular.
  \item \textbf{HTML / CSS} : structuration et mise en forme des interfaces.
\end{itemize}

\subsection{Architecture applicative et communication}
Les échanges entre le front-end et le back-end s'effectuent via des API REST exposées par la
passerelle. La communication asynchrone repose sur RabbitMQ : lorsqu'un événement métier se
produit (inscription, confirmation, etc.), un message est publié sur un \textit{exchange} et
routé vers une file consommée par le service de notification. L'extrait suivant présente la
configuration de la file de messages.

\begin{lstlisting}[language=Java,caption={Configuration RabbitMQ du service de notification},label={lst:rabbit}]
@Configuration
public class RabbitMQConfig {
    public static final String QUEUE_NAME = "notification.queue";
    public static final String EXCHANGE_NAME = "event.exchange";
    public static final String ROUTING_KEY = "notification.routing.key";

    @Bean public Queue queue() { return new Queue(QUEUE_NAME, true); }
    @Bean public TopicExchange exchange() { return new TopicExchange(EXCHANGE_NAME); }
    @Bean public Binding binding(Queue queue, TopicExchange exchange) {
        return BindingBuilder.bind(queue).to(exchange).with(ROUTING_KEY);
    }
}
\end{lstlisting}

\subsection{Aperçu de l'environnement technique}
\begin{figure}[H]\centering
  \ph[0.85\textwidth]{[ Planche des logos des technologies : Angular, Spring Boot, Spring Cloud, Keycloak, MongoDB, PostgreSQL, RabbitMQ, MailHog, Docker, OpenShift, Maven, Git ]}
  \caption{Ensemble des technologies et outils mobilisés pour le développement de la plateforme}\label{fig:env-tech}
\end{figure}

\begin{figure}[H]\centering
  \ph[0.85\textwidth]{[ Console web Eureka listant les instances UP : api-gateway, user, event, registration, notification, ticket, feedback, ai ]}
  \caption{Tableau de bord Eureka affichant les microservices enregistrés}\label{fig:env-eureka}
\end{figure}

\begin{figure}[H]\centering
  \ph[0.85\textwidth]{[ Console d'administration RabbitMQ : files de notifications, exchanges et connexions actifs ]}
  \caption{Interface de gestion RabbitMQ (messagerie asynchrone)}\label{fig:env-rabbit}
\end{figure}

\begin{figure}[H]\centering
  \ph[0.85\textwidth]{[ Interface web MailHog : liste des courriels de notification interceptés (objet, expéditeur, aperçu) ]}
  \caption{Boîte de réception MailHog interceptant les courriels de la plateforme}\label{fig:env-mailhog}
\end{figure}

\section{Sprint 1 : Authentification et gestion des utilisateurs}

\subsection{Objectif du sprint}
Ce premier sprint met en place les fondations de la plateforme : l'authentification
centralisée via Keycloak, la résolution des rôles et la gestion des collaborateurs par
l'administrateur.

\subsection{Sprint Backlog}
\begin{table}[H]\centering\renewcommand{\arraystretch}{1.25}
\begin{tabular}{|p{1cm}|p{10.5cm}|p{1.6cm}|}\hline
\rowcolor{gray!15}\textbf{ID} & \textbf{Tâche} & \textbf{Statut} \\ \hline
T1 & Configurer le realm Keycloak et les clients OAuth2. & Terminé \\ \hline
T2 & Intégrer l'authentification côté front-end (keycloak-angular). & Terminé \\ \hline
T3 & Résoudre le rôle effectif par priorité (Admin > Organisateur > Participant). & Terminé \\ \hline
T4 & Développer la gestion des utilisateurs (CRUD) côté administrateur. & Terminé \\ \hline
T5 & Implémenter l'attribution et la modification des rôles. & Terminé \\ \hline
\end{tabular}
\caption{Sprint Backlog du Sprint 1}\label{tab:sb1}
\end{table}

\subsection{Réalisation}
L'authentification et la gestion des identités reposent sur \textbf{Keycloak}, un serveur
de gestion des identités et des accès (IAM) que nous avons intégré comme fournisseur
d'identité central de la plateforme.

\subsubsection{Configuration du serveur Keycloak}
Nous avons configuré un \textit{realm} dédié, \texttt{pfe-events}, qui regroupe l'ensemble des
utilisateurs, des rôles et des clients de l'application. Deux clients y sont déclarés :

\begin{table}[H]\centering\renewcommand{\arraystretch}{1.3}
\begin{tabular}{|p{3.8cm}|p{2.6cm}|p{7cm}|}\hline
\rowcolor{gray!15}\textbf{Client} & \textbf{Type} & \textbf{Rôle} \\ \hline
\texttt{pfe-events-frontend} & Public (SPA) & Authentifie l'application Angular via le flux
\textit{Authorization Code} avec PKCE. \\ \hline
\texttt{pfe-events-api} & Confidentiel & Représente les services back-end (serveurs de
ressources) qui valident les jetons. \\ \hline
\end{tabular}
\caption{Clients Keycloak du realm pfe-events}\label{tab:kc-clients}
\end{table}

Trois rôles applicatifs sont définis : \texttt{ADMIN}, \texttt{ORGANISATEUR} et
\texttt{PARTICIPANT}. Le realm attribuant par défaut le rôle \textit{Participant} à tout
nouvel utilisateur, le front-end résout le \textbf{rôle effectif} par ordre de priorité
(\textsc{Admin} > \textsc{Organisateur} > \textsc{Participant}), afin de garantir une stricte
séparation des vues.

\subsubsection{Flux d'authentification (OAuth2 / OpenID Connect)}
Le processus s'appuie sur les standards \textbf{OAuth2} et \textbf{OpenID Connect} :
\begin{enumerate}[leftmargin=1.2cm]
  \item l'utilisateur est redirigé vers la page de connexion Keycloak (thème LINSOFT) ;
  \item après authentification, Keycloak délivre un \textbf{jeton JWT} signé ;
  \item ce jeton accompagne chaque requête ; l'\textbf{API Gateway} en valide la signature et
  l'expiration, puis convertit les rôles du realm en autorités Spring ;
  \item chaque microservice re-valide le jeton (défense en profondeur).
\end{enumerate}
À la première connexion, le compte du collaborateur est automatiquement \textbf{provisionné}
dans le service utilisateur.

\subsubsection{Personnalisation de la page de connexion}
Un \textbf{thème de connexion personnalisé} aux couleurs de LINSOFT (logo, arrière-plan,
charte graphique rouge) a été développé et intégré à Keycloak, afin d'offrir une expérience de
connexion cohérente avec l'identité de l'entreprise.

\begin{figure}[H]\centering
  \ph[0.72\textwidth]{[ Capture : page de connexion Keycloak (thème LINSOFT) ]}
  \caption{Page de connexion personnalisée LINSOFT}\label{fig:kc-login}
\end{figure}

\begin{figure}[H]\centering
  \ph[0.85\textwidth]{[ Capture : console d'administration Keycloak (realm, rôles, clients) ]}
  \caption{Console d'administration Keycloak}\label{fig:kc-admin}
\end{figure}

\subsubsection{Gestion des collaborateurs}
L'administrateur dispose d'une interface de gestion des collaborateurs (consultation,
création, modification, suppression) et peut y attribuer ou modifier les rôles.

\begin{figure}[H]\centering
  \ph[0.9\textwidth]{[ Capture : gestion des collaborateurs (administrateur) ]}
  \caption{Interface de gestion des collaborateurs}\label{fig:users}
\end{figure}

\begin{figure}[H]\centering
  \ph[0.85\textwidth]{[ Onglet Roles du realm pfe-events : rôles Participant, Organisateur, Administrateur ]}
  \caption{Définition des rôles applicatifs au niveau du realm Keycloak}\label{fig:kc-roles}
\end{figure}

\begin{figure}[H]\centering
  \ph[0.85\textwidth]{[ Formulaire « Add user » de Keycloak : username, email, prénom, nom, activation ]}
  \caption{Création d'un utilisateur dans la console d'administration Keycloak}\label{fig:kc-adduser}
\end{figure}

\begin{figure}[H]\centering
  \ph[0.85\textwidth]{[ Menu d'affectation de rôle (Participant / Organisateur / Administrateur) sur une ligne collaborateur ]}
  \caption{Attribution et modification du rôle d'un collaborateur}\label{fig:role-assign}
\end{figure}

\section{Sprint 2 : Gestion des sessions et des inscriptions}

\subsection{Objectif du sprint}
Ce sprint couvre le cœur métier de la plateforme. Il permet d'abord aux organisateurs de créer
et de gérer leurs sessions, et à l'administrateur de les modérer avant publication au
catalogue ; il implémente ensuite la demande d'inscription par le participant et son workflow
de validation, notifications comprises.

\subsection{Sprint Backlog}
\begin{table}[H]\centering\renewcommand{\arraystretch}{1.25}
\begin{tabular}{|p{1cm}|p{10.5cm}|p{1.6cm}|}\hline
\rowcolor{gray!15}\textbf{ID} & \textbf{Tâche} & \textbf{Statut} \\ \hline
T6 & Développer le formulaire de création de session. & Terminé \\ \hline
T7 & Gérer les types, modes et catégories de sessions. & Terminé \\ \hline
T8 & Développer la modération des sessions (approuver / refuser). & Terminé \\ \hline
T9 & Afficher le catalogue des sessions publiées. & Terminé \\ \hline
T10 & Développer la demande d'inscription (statut PENDING). & Terminé \\ \hline
T11 & Développer la validation / le refus des inscriptions. & Terminé \\ \hline
T12 & Gérer la décrémentation des places disponibles. & Terminé \\ \hline
T13 & Développer les onglets de suivi des statuts d'inscription. & Terminé \\ \hline
\end{tabular}
\caption{Sprint Backlog du Sprint 2}\label{tab:sb2}
\end{table}

\subsection{Gestion du catalogue des sessions}
\begin{figure}[H]\centering
  \ph[0.9\textwidth]{[ Capture : formulaire de création d'une session ]}
  \caption{Formulaire de création d'une session}\label{fig:create-event}
\end{figure}

\begin{figure}[H]\centering
  \ph[0.9\textwidth]{[ Capture : catalogue des sessions ]}
  \caption{Catalogue des sessions}\label{fig:catalogue}
\end{figure}

\begin{figure}[H]\centering
  \ph[0.85\textwidth]{[ Page « Mes sessions » : sessions créées, badges de statut (En attente / Approuvée / Refusée), actions Éditer et Supprimer ]}
  \caption{Liste « Mes sessions » de l'organisateur avec les statuts de modération}\label{fig:my-sessions}
\end{figure}

\begin{figure}[H]\centering
  \ph[0.85\textwidth]{[ Écran d'édition : formulaire pré-rempli (titre, type, mode, places, catégorie, niveau) prêt à être modifié ]}
  \caption{Édition d'une session existante par l'organisateur}\label{fig:edit-session}
\end{figure}

\begin{figure}[H]\centering
  \ph[0.85\textwidth]{[ Page de modération : sessions en attente avec leurs détails et les boutons Approuver et Refuser ]}
  \caption{Interface de modération des sessions (administrateur)}\label{fig:moderation}
\end{figure}

\begin{figure}[H]\centering
  \ph[0.85\textwidth]{[ Détail d'une session : description, type, mode, catégorie, niveau, certification, places, bouton d'inscription ]}
  \caption{Page de détail d'une session consultée par le participant}\label{fig:session-detail}
\end{figure}

\subsection{Inscriptions et workflow de validation}
\begin{figure}[H]\centering
  \ph[0.9\textwidth]{[ Capture : liste des inscriptions à valider (administrateur) ]}
  \caption{Validation des inscriptions}\label{fig:validation}
\end{figure}

\begin{figure}[H]\centering
  \ph[0.9\textwidth]{[ Capture : mes inscriptions et leurs statuts (participant) ]}
  \caption{Suivi des inscriptions du participant}\label{fig:my-reg}
\end{figure}

\begin{figure}[H]\centering
  \ph[0.85\textwidth]{[ Détail d'une session avec le bouton rouge LINSOFT « Demander une inscription » ]}
  \caption{Demande d'inscription depuis le détail d'une session}\label{fig:demande-insc}
\end{figure}

\begin{figure}[H]\centering
  \ph[0.85\textwidth]{[ Vue planning / calendrier des sessions disponibles côté participant ]}
  \caption{Vue planning des sessions accessible au participant}\label{fig:planning}
\end{figure}

\begin{figure}[H]\centering
  \ph[0.85\textwidth]{[ Sous-onglets « Approuvées » et « Refusées » : historique des demandes traitées (participant, session) ]}
  \caption{Traçabilité des demandes traitées (sous-onglets Approuvées / Refusées)}\label{fig:val-tabs}
\end{figure}

\begin{figure}[H]\centering
  \ph[0.85\textwidth]{[ Notification in-app affichée au participant après validation ou refus de sa demande ]}
  \caption{Notification in-app indiquant le nouveau statut de l'inscription}\label{fig:notif-inapp}
\end{figure}

\section*{Conclusion}
Cette première release a livré le cœur fonctionnel de la plateforme : authentification et
gestion des utilisateurs, gestion du catalogue et workflow d'inscription. Le chapitre suivant
présente la seconde release, dédiée aux fonctionnalités à valeur ajoutée, à l'industrialisation
et au déploiement.
\clearpage